\section{Discussion}
\label{sec:discussion}

\para{Symbolic composition works after perception is grounded.} The experiment supports the central neural-symbolic hypothesis in this controlled setting. When the digit CNN learns accurate predicates, exact marginalization over the addition table produces valid sums outside the supervised label support. This is why \symbolicsum can reach 80.6\% accuracy on sums 14--18 even though those labels never appear during training. The low-budget condition also clarifies the boundary of the claim: symbolic structure cannot compensate for poorly grounded predicates.

\para{The simplest symbolic likelihood is the strongest model.} We expected the anchors and posterior entropy to stabilize latent digit assignments. They do help relative to the direct baseline at 5,000 labels, but they do not improve over \symbolicsum. One explanation is that the sum likelihood already provides enough signal once 5,000 pair labels are available, while the anchor and entropy terms add constraints that are too strong early in training. A better variant may need a schedule that delays posterior sharpening or reduces anchor weight after the digit classes become identifiable.

\para{The direct baseline fails for the intended reason.} \directsum is not a universal neural baseline; it is a direct classifier under a deliberately harsh label-support shift. It can improve on IID restricted pairs, but the missing labels 14--18 receive no positive supervised examples. The point of the comparison is therefore not that direct CNNs cannot learn \mnist arithmetic in general. The point is that a direct label map and a symbolic composition model make different extrapolation commitments when the label support changes.

\para{Limitations.} \addmnist is synthetic and much simpler than language reasoning, visual question answering, or knowledge-graph reasoning. We use an exact known rule, so the experiment does not test rule induction. The OOD split changes digit-pair composition but still uses the same \mnist image distribution, so it is not a visual domain-shift benchmark. We also run only three seeds. The OOD effects at 5,000 labels are large and consistent, but more seeds and broader hyperparameter sweeps would give a stronger statistical picture.

\para{Broader implications.} The results argue for treating exact symbolic marginalization as a strong baseline whenever the rule is known and the latent predicate space is small. It is easy to overlook this baseline in favor of heavier neural-symbolic frameworks, but this experiment shows that the simple version can already deliver the desired label extrapolation. The next challenge is to preserve this advantage as the predicate space grows, the rule becomes approximate, or perception shifts away from the training domain.
